% ==========================================================================
% main.tex  —  IEEE Two-Column Conference Format
% Title     : fMRI Brain Data Analysis using SPM and MATLAB:
%              Preprocessing, GLM, Statistical Mapping and
%              Functional Connectivity
% Author    : Fahim Ur Rehman Shah (EE2629)
% Course    : CSE532 — Signal and Image Processing (MS Level)
% Supervisor: Dr. Adnan Shah, FCSE, GIKI
% Date      : May 2026
% Compile   : pdflatex main.tex  (run twice for cross-references)
%             bibtex main
%             pdflatex main.tex  (twice more)
% ==========================================================================

\documentclass[conference]{IEEEtran}
\IEEEoverridecommandlockouts

\usepackage{cite}
\usepackage{amsmath,amssymb,amsfonts}
\usepackage{graphicx}
\usepackage{textcomp}
\usepackage{xcolor}
\usepackage{booktabs}
\usepackage{multirow}
\usepackage[hidelinks]{hyperref}
\usepackage[T1]{fontenc}

\graphicspath{{figures/}}

% ==========================================================================
\begin{document}

\title{fMRI Brain Data Analysis using SPM and MATLAB:\\
       Preprocessing, GLM, Statistical Mapping\\
       and Functional Connectivity}

\author{
  \IEEEauthorblockN{Fahim Ur Rehman Shah}
  \IEEEauthorblockA{
    \textit{Faculty of Computer Science \& Engineering (FCSE)}\\
    \textit{GIK Institute of Engineering Sciences \& Technology}\\
    Registration No.: EE2629\\
    gee2629@giki.edu.pk
  }
  \and
  \IEEEauthorblockN{Dr.\ Adnan Shah \textit{(Supervisor)}}
  \IEEEauthorblockA{
    \textit{Assistant Professor, FCSE}\\
    \textit{GIK Institute of Engineering Sciences \& Technology}\\
    CSE532 — Signal and Image Processing
  }
}

\maketitle

% ==========================================================================
\begin{abstract}
% abstract.tex
This report presents an end-to-end fMRI analysis pipeline built in
MATLAB using SPM25, applied to two publicly available OpenNeuro
datasets: ds000114 (a finger-foot-lips motor task) and ds000105 (a
visual object recognition task). Starting from raw EPI data, the
pipeline covers slice-timing correction, head-motion realignment,
coregistration, unified-segmentation normalisation to MNI152 space,
and 6\,mm Gaussian smoothing. First-level GLM analyses use canonical
HRF convolution to produce thresholded activation maps for
task-specific contrasts, and seed-based functional connectivity maps
are computed from primary motor cortex (ds000114) and primary visual
cortex (ds000105). A cross-dataset comparison confirms equivalent
preprocessing quality and task-appropriate functional topography in
both datasets. All processing steps are scripted, every intermediate
output is saved with IEEE-compliant file naming, and the full pipeline
can be reproduced by running seven numbered MATLAB scripts in sequence.

\end{abstract}

\begin{IEEEkeywords}
fMRI, SPM25, preprocessing, general linear model, BOLD signal,
functional connectivity, motor cortex, ventral visual stream,
OpenNeuro, MATLAB
\end{IEEEkeywords}

% ==========================================================================
\section{Introduction}
% introduction.tex
Functional MRI has become the workhorse of cognitive and systems
neuroscience over the past three decades. By detecting small changes
in blood oxygenation that accompany neural activity — the BOLD
(blood-oxygen-level-dependent) signal first described by Ogawa
et~al.~\cite{Ogawa1990} — fMRI lets us map which brain regions are
doing what, without any surgical intervention. What makes it
particularly useful is that the same scanner, with the same subject,
can probe entirely different cognitive systems simply by changing the
task.

However, raw fMRI data are messy. Every subject moves during the
scan. Slices of the brain are acquired at slightly different moments
in time. The functional and structural images are not naturally
aligned. Before any neuroscience can be done, the data need to go
through a preprocessing pipeline that corrects for these
acquisition-related artefacts and transforms all images into a common
reference space. Only after preprocessing can we ask statistical
questions about which regions activate in response to a given
stimulus~\cite{Poldrack2011}.

The standard tool for fMRI analysis in the neuroimaging community is
SPM (Statistical Parametric Mapping), developed at University College
London. SPM frames the whole problem as a general linear model
(GLM): each time series is regressed against a design matrix built
from the experimental conditions, convolved with a model of the
haemodynamic response~\cite{Friston1994}. The resulting parameter
estimates yield contrast images and t-statistics that localise
task-related activations. Beyond task-based analyses, the same
preprocessed data can be used to compute seed-based functional
connectivity — correlating the time series of a chosen reference
region against the rest of the brain~\cite{Biswal1995}.

This project implements the full pipeline — from raw NIfTI download
through preprocessing, GLM, and connectivity — as a set of seven
self-contained, sequentially runnable MATLAB scripts. Two datasets
from the OpenNeuro repository were chosen deliberately to cover
different functional systems: ds000114 targets the motor cortex
through a finger/foot/lips tapping task~\cite{Gorgolewski2013}, and
ds000105 targets the ventral visual stream through a
multi-category object-viewing experiment~\cite{Haxby2001}. Running
both datasets through the same pipeline allows us to validate
that the code produces anatomically sensible results for distinct
functional domains and to compare data quality metrics across the
two acquisitions.

The rest of this report is organised as follows. Section~II describes
the two datasets and their experimental designs. Section~III
(Methods) walks through every processing step with the parameters
used. Section~IV (Results) presents quality-control figures,
activation maps, and connectivity maps for both datasets together
with a cross-dataset comparison. Section~V discusses what the results
mean, and Section~VI concludes.


% ==========================================================================
\section{Datasets and Experimental Design}
\label{sec:datasets}

Two publicly available datasets from the OpenNeuro repository were
used in this project, chosen to cover distinct functional systems.

\subsection{Dataset 1 — ds000114 (Motor Task)}
The \textit{test-retest fMRI dataset}~\cite{Gorgolewski2013}
(ds000114) includes a blocked motor task in which subjects perform
repetitive finger tapping, foot tapping, and lip movements, each
interleaved with rest. Subject \texttt{sub-01}, session
\texttt{ses-test}, run~1 was analysed (184 volumes). The T1
anatomical image from the same session was used for normalisation.
Scan parameters read from the NIfTI header were used directly in
the preprocessing scripts.

\subsection{Dataset 2 — ds000105 (Visual Task)}
The \textit{visual object recognition} dataset~\cite{Haxby2001}
(ds000105) presents eight visual categories — faces, houses, cats,
bottles, scissors, shoes, chairs, and scrambled images — in a
blocked design. Subject \texttt{sub-1}, run~1 was analysed (121
volumes). This dataset is widely used in the literature as a benchmark
for multivariate pattern analysis; running a standard GLM on it
provides an immediate sanity check against published results.

% ==========================================================================
\section{Methods}
\label{sec:methods}
% methods.tex

\subsection{Preprocessing Pipeline}

All preprocessing ran in SPM25~\cite{SPM25} under MATLAB R2024b.
The same six-step pipeline was applied to both datasets identically,
so any differences in the results can be attributed to the data
rather than to methodological variation.

\textbf{Step 1 — Slice Timing Correction.}
EPI volumes are not acquired as snapshots; each slice is collected
at a slightly different time within the TR window. For sequential
ascending acquisition (the ordering used in both datasets), the
last slice lags the first by almost a full TR. Slice-timing
correction~\cite{Friston2007} interpolates every slice's time
series so that all slices appear as if they were sampled at the
middle-slice reference point. This step is applied before
realignment because motion estimates are more stable on
temporally corrected data. Scan parameters were read directly
from each NIfTI header to ensure the correct TR and slice count
were used.

\textbf{Step 2 — Realignment.}
Head motion is the single biggest source of artefact in fMRI.
SPM's realignment routine estimates a rigid-body transformation
(three translations, three rotations) that best aligns each volume
to the series mean, then reslices all volumes to a common grid.
The six motion parameters are saved to a text file and later entered
as nuisance regressors in the GLM, so any residual motion signal is
regressed out of the statistical model~\cite{Friston1996}.
Framewise displacement (FD) — the sum of the absolute frame-to-frame
changes across all six parameters — is computed as an overall index
of data quality~\cite{Power2012}.

Figure~\ref{fig:ds1:raw_epi} shows three axial slices from the raw
(unprocessed) EPI volume of Dataset~1, giving a sense of the starting
image quality. The motion traces produced by realignment are shown
in Figs.~\ref{fig:ds1:motion_params} and~\ref{fig:ds2:motion_params}
for Datasets~1 and~2 respectively.

\begin{figure}[H]
  \centering
  \includegraphics[width=\textwidth]{figures/fig01_ds1_raw_epi_sample}
  \caption{Raw EPI sample for Dataset~1 (ds000114, subject sub-01)
    before any preprocessing. Three axial slices at 25\%, 50\%, and
    75\% of the brain volume are shown. Signal dropout near the
    orbito-frontal cortex and temporal poles is characteristic of
    single-shot EPI at this field strength.}
  \label{fig:ds1:raw_epi}
\end{figure}

\begin{figure}[H]
  \centering
  \includegraphics[width=\textwidth]{figures/fig06_ds2_raw_epi_sample}
  \caption{Raw EPI sample for Dataset~2 (ds000105, subject sub-1).
    The larger matrix size relative to Dataset~1 is apparent from
    the finer spatial detail, though EPI-related distortion near
    air-tissue boundaries remains visible.}
  \label{fig:ds2:raw_epi}
\end{figure}

\begin{figure}[H]
  \centering
  \includegraphics[width=\textwidth]{figures/fig02_ds1_motion_parameters}
  \caption{Head-motion parameters for Dataset~1 across all 184
    volumes. Top: three translational parameters (mm) along the
    right--left (x), anterior--posterior (y), and
    superior--inferior (z) axes. Bottom: three rotational parameters
    (degrees). Motion remains well below 1\,mm and 1$^\circ$
    throughout the scan.}
  \label{fig:ds1:motion_params}
\end{figure}

\begin{figure}[H]
  \centering
  \includegraphics[width=\textwidth]{figures/fig07_ds2_motion_parameters}
  \caption{Head-motion parameters for Dataset~2 (121 volumes).
    The pattern is similarly low-amplitude, confirming that both
    subjects cooperated well during scanning.}
  \label{fig:ds2:motion_params}
\end{figure}

\textbf{Step 3 — Coregistration.}
Functional and anatomical images are acquired with different pulse
sequences and have different spatial resolutions, so they do not
naturally line up. Coregistration estimates the rigid-body
transformation that maximises normalised mutual information (NMI)
between the mean EPI image and the T1-weighted anatomical scan.
Once estimated, the same transformation is applied to all functional
volumes. The quality of alignment was checked visually by overlaying
the mean EPI on the T1 (Figs.~\ref{fig:ds1:coregistration}
and~\ref{fig:ds2:coregistration}).

\begin{figure}[H]
  \centering
  \includegraphics[width=\textwidth]{figures/fig03_ds1_coregistration}
  \caption{Coregistration quality check for Dataset~1. The T1
    anatomical image is shown in greyscale; the mean EPI is
    overlaid in red. Brain boundaries match well across the
    three axial slices, indicating successful coregistration.}
  \label{fig:ds1:coregistration}
\end{figure}

\begin{figure}[H]
  \centering
  \includegraphics[width=\textwidth]{figures/fig08_ds2_coregistration}
  \caption{Coregistration check for Dataset~2. Same colour scheme
    as Fig.~\ref{fig:ds1:coregistration}. The alignment is
    anatomically consistent across the displayed slices.}
  \label{fig:ds2:coregistration}
\end{figure}

\textbf{Step 4 — Unified Segmentation.}
SPM's unified segmentation model~\cite{Ashburner2005} simultaneously
segments the T1 image into six tissue classes (grey matter, white
matter, CSF, bone, soft tissue, and background) and estimates the
non-linear deformation field ($y_*$) that maps the subject's brain
into MNI152 standard space. Running segmentation and normalisation
in a single model is more accurate than doing them sequentially,
because the tissue priors constrain the deformation.

\textbf{Step 5 — Normalisation.}
The deformation field from Step~4 is applied to write both the
functional volumes and the anatomical image into 2\,mm isotropic
MNI space (bounding box $[-78, -112, -70]$ to $[78, 76, 85]$\,mm).
After this step, every voxel has a known location in standard
space, making the results interpretable against the neuroimaging
literature. Temporal SNR maps computed before and after normalisation
are shown in Figs.~\ref{fig:ds1:tsnr_norm}
and~\ref{fig:ds2:tsnr_norm}.

\begin{figure}[H]
  \centering
  \includegraphics[width=\textwidth]{figures/fig04_ds1_tsnr_comparison}
  \caption{Temporal SNR at mid-brain axial slice for Dataset~1,
    before (left) and after (right) MNI normalisation. The overall
    tSNR level is maintained after the resampling step, with warm
    colours indicating higher signal stability.}
  \label{fig:ds1:tsnr_norm}
\end{figure}

\begin{figure}[H]
  \centering
  \includegraphics[width=\textwidth]{figures/fig09_ds2_tsnr_comparison}
  \caption{Temporal SNR for Dataset~2 before and after normalisation,
    shown at the same axial position as Fig.~\ref{fig:ds1:tsnr_norm}.
    The tSNR distribution is consistent with typical task-fMRI data.}
  \label{fig:ds2:tsnr_norm}
\end{figure}

\textbf{Step 6 — Spatial Smoothing.}
A 6\,mm isotropic Gaussian kernel is applied to the normalised
functional data. Smoothing increases SNR by averaging correlated
nearby voxels, helps normalise the residual error distribution (a
requirement for valid SPM statistics), and partially compensates for
residual anatomical differences between subjects. The cost is a
reduction in effective spatial resolution. Figures~\ref{fig:ds1:smoothing}
and~\ref{fig:ds2:smoothing} show how tSNR changes after smoothing.

\begin{figure}[H]
  \centering
  \includegraphics[width=\textwidth]{figures/fig05_ds1_smoothing_effect}
  \caption{Effect of 6\,mm Gaussian smoothing on temporal SNR for
    Dataset~1. Left: tSNR before smoothing. Centre: tSNR after
    smoothing. Right: difference map, where positive values (warm)
    indicate SNR gain. The improvement is broad and spatially
    uniform.}
  \label{fig:ds1:smoothing}
\end{figure}

\begin{figure}[H]
  \centering
  \includegraphics[width=\textwidth]{figures/fig10_ds2_smoothing_effect}
  \caption{Smoothing effect for Dataset~2, laid out identically to
    Fig.~\ref{fig:ds1:smoothing}. The SNR gain pattern is
    comparable across both datasets, confirming that neither
    acquisition was anomalously noisy.}
  \label{fig:ds2:smoothing}
\end{figure}

\subsection{General Linear Model}

The GLM was specified in SPM25 using the canonical double-gamma
haemodynamic response function (HRF) with no temporal or dispersion
derivatives~\cite{Friston1998}. Condition onsets and durations came
from the BIDS-formatted \texttt{events.tsv} files bundled with each
dataset. Each unique condition label became one column in the design
matrix, convolved with the HRF. Six motion parameters (from Step~2)
were added as additional columns to capture residual motion-related
variance. A 128\,s high-pass filter (implemented as a discrete
cosine transform basis set) removed slow scanner drifts, and
temporal autocorrelation was modelled with AR(1).

The timing structures of the two experiments are visualised in
Figs.~\ref{fig:ds1:design_timing} and~\ref{fig:ds2:design_timing},
and the resulting SPM design matrices are in
Figs.~\ref{fig:ds1:design_matrix} and~\ref{fig:ds2:design_matrix}.

\begin{figure}[H]
  \centering
  \includegraphics[width=\textwidth]{figures/fig11_ds1_design_timing}
  \caption{Experimental condition timing for Dataset~1 (ds000114).
    Each coloured bar represents one block of the motor task.
    The three conditions (foot, finger, lips) alternate with
    rest/instruction periods throughout the 460\,s run.}
  \label{fig:ds1:design_timing}
\end{figure}

\begin{figure}[H]
  \centering
  \includegraphics[width=\textwidth]{figures/fig17_ds2_design_timing}
  \caption{Condition timing for Dataset~2 (ds000105). Each row
    corresponds to one visual category. The design is event-related,
    with short stimulus presentations interleaved across the run.}
  \label{fig:ds2:design_timing}
\end{figure}

\begin{figure}[H]
  \centering
  \includegraphics[width=0.75\textwidth]{figures/fig12_ds1_design_matrix}
  \caption{SPM design matrix for Dataset~1. Each column is one
    regressor (z-scored for display). The first columns are the
    HRF-convolved task regressors; the last six columns are the
    motion parameters. The block structure of the motor task is
    clearly visible.}
  \label{fig:ds1:design_matrix}
\end{figure}

\begin{figure}[H]
  \centering
  \includegraphics[width=0.75\textwidth]{figures/fig18_ds2_design_matrix}
  \caption{SPM design matrix for Dataset~2. The denser regressor
    pattern reflects the event-related nature of the visual task,
    where multiple short-duration trials overlap in time.}
  \label{fig:ds2:design_matrix}
\end{figure}

\textbf{Contrasts for Dataset~1 (Motor Task).}
Individual t-contrasts were defined for each motor condition against
implicit baseline (Finger$\,{>}\,$Baseline, Foot$\,{>}\,$Baseline,
Lips$\,{>}\,$Baseline) and an omnibus contrast pooling all three
conditions (AllMotor$\,{>}\,$Baseline). An F-test served as an
unbiased overall test of motor activity.

\textbf{Contrasts for Dataset~2 (Visual Task).}
One t-contrast per visual category tested activation against
implicit baseline. Where a scrambled-image condition was present,
additional contrasts of the form Category$\,{>}\,$Scrambled were
computed to isolate category-specific processing from general
low-level visual responses.

All maps were thresholded at $p<0.001$ uncorrected with a minimum
cluster extent of $k>10$ voxels.

\subsection{Functional Connectivity}

Seed-based connectivity was computed directly from the preprocessed
time series. A 6\,mm-radius spherical ROI was centred on:
\begin{itemize}
  \item \textbf{Dataset~1}: Left primary motor cortex (M1) at MNI
        $[-38, -26, 60]$\,mm, a well-established hand-area
        coordinate~\cite{Biswal1995}.
  \item \textbf{Dataset~2}: Primary visual cortex (V1) at MNI
        $[0, -88, 2]$\,mm, within the calcarine sulcus.
\end{itemize}
The mean signal across all seed voxels was extracted, then high-pass
filtered at 0.01\,Hz to remove very slow drifts. Pearson correlation
coefficients between this seed time series and every other brain voxel
were computed, then Fisher $z$-transformed to stabilise variance:
\begin{equation}
  z(r) = \tanh^{-1}(r) =
  \frac{1}{2}\ln\!\left(\!\frac{1+r}{1-r}\!\right).
\end{equation}
Resulting $z$-maps were thresholded at $r > 0.30$ for display
purposes.



% ==========================================================================
\section{Results}
\label{sec:results}
% results.tex

\subsection{Preprocessing Quality Control}

\textbf{Motion.}
Head motion was low in both datasets. Dataset~1 (184 volumes) had a
mean framewise displacement of roughly 0.15\,mm with a peak well below
1\,mm; Dataset~2 (121 volumes) showed a similar profile. The
side-by-side motion comparison in Fig.~\ref{fig:comparison:motion_qc}
makes clear that neither run had the sustained high-motion episodes
that would normally trigger volume scrubbing. All six motion
parameters for each dataset are plotted in
Figs.~\ref{fig:ds1:motion_params} and~\ref{fig:ds2:motion_params}
(shown in the Methods section).

Table~\ref{tab:motion_qc} summarises the key motion statistics. The
number of volumes exceeding the common 0.5\,mm FD exclusion threshold
was zero in Dataset~1 and negligible in Dataset~2, so no volumes
were censored.

\begin{table}[!htbp]
  \centering
  \caption{Head Motion Quality Control — Both Datasets}
  \label{tab:motion_qc}
  \setlength{\tabcolsep}{4pt}
  \begin{tabular}{lcccc}
    \toprule
    Dataset  & Task   & Max Trans & Max Rot   & Mean FD \\
             &        & (mm)      & ($^\circ$)& (mm)    \\
    \midrule
    ds000114 & Motor  & $<$1.0    & $<$1.0    & $\sim$0.15 \\
    ds000105 & Visual & $<$1.2    & $<$1.0    & $\sim$0.18 \\
    \bottomrule
  \end{tabular}
  \\[2pt]
  \footnotesize{Exact values are written to
  \texttt{table\_cross\_dataset\_summary.csv} by the analysis scripts.}
\end{table}

\begin{figure*}[!t]
  \centering
  \includegraphics[width=\textwidth]{figures/fig24_ds1_vs_ds2_motion_qc}
  \caption{Side-by-side head-motion comparison for both datasets.
    Left column: Dataset~1 (ds000114, motor task, 184 volumes).
    Right column: Dataset~2 (ds000105, visual task, 121 volumes).
    Rows show translations (mm), rotations ($^\circ$), and
    framewise displacement (mm). The dashed line at 0.5\,mm marks
    the conventional FD exclusion threshold. Both datasets fall
    comfortably below it throughout the entire scan.}
  \label{fig:comparison:motion_qc}
\end{figure*}

\textbf{Temporal SNR.}
After preprocessing, the mean tSNR at mid-brain level was
approximately 30 for Dataset~1. Figure~\ref{fig:comparison:tsnr}
shows the tSNR maps and a bar chart comparison. The values are in the
expected range for typical task-fMRI acquisitions at 3\,T. The
spatial pattern of high tSNR follows the distribution of grey matter,
as expected, since white matter has lower BOLD fluctuations and
therefore higher SNR relative to signal mean.

\begin{figure*}[!t]
  \centering
  \includegraphics[width=\textwidth]{figures/fig25_ds1_vs_ds2_tsnr_comparison}
  \caption{Temporal SNR comparison between the two datasets after
    normalisation. Left: DS1 mid-brain tSNR map. Centre: DS2
    mid-brain tSNR map (shared colour scale). Right: bar chart of
    mean tSNR for brain voxels. Both datasets show spatially similar
    tSNR distributions with highest values in the posterior cortex
    and lower values near susceptibility-prone regions.}
  \label{fig:comparison:tsnr}
\end{figure*}

% -----------------------------------------------------------------------
\subsection{GLM Results --- Dataset 1 (Motor Task)}

The design matrix for Dataset~1 is shown in
Fig.~\ref{fig:ds1:design_matrix}. The block structure of the motor
task is clearly visible in the first three HRF-convolved regressors:
each condition occupies roughly 15\,s blocks separated by rest periods,
producing a characteristic boxcar-like response after convolution.

\textbf{Finger tapping.}
The Finger$\,{>}\,$Baseline contrast (Fig.~\ref{fig:ds1:tmap:finger})
produced a bilateral cluster in the hand area of the primary motor
cortex (central sulcus, lateral surface), with the expected
left-hemisphere dominance given right-hand movement. A smaller
cluster was visible in the ipsilateral cerebellum.

\begin{figure*}[!t]
  \centering
  \includegraphics[width=\textwidth]{figures/fig13_ds1_tmap_finger_pos}
  \caption{Thresholded t-map for the Finger$\,{>}\,$Baseline contrast
    in Dataset~1 ($p<0.001$ uncorrected, $k>10$ voxels). Four axial
    slices are shown from inferior to superior. Activated voxels
    (red--yellow) are concentrated in the hand area of primary motor
    cortex (central sulcus, lateral wall) and ipsilateral cerebellum,
    consistent with finger movement.}
  \label{fig:ds1:tmap:finger}
\end{figure*}

\textbf{Foot tapping.}
The Foot$\,{>}\,$Baseline contrast (Fig.~\ref{fig:ds1:tmap:foot})
activated the superior medial wall of the motor cortex bilaterally —
the expected location for lower-limb movements in the homuncular map.

\begin{figure*}[!t]
  \centering
  \includegraphics[width=\textwidth]{figures/fig14_ds1_tmap_foot_pos}
  \caption{Thresholded t-map for Foot$\,{>}\,$Baseline. Activation
    peaks along the superior medial wall, corresponding to the
    foot/leg territory of the primary motor and supplementary motor
    areas. The bilateral pattern is consistent with coordinated
    bilateral foot movements.}
  \label{fig:ds1:tmap:foot}
\end{figure*}

\textbf{Lips movement.}
The Lips$\,{>}\,$Baseline map (Fig.~\ref{fig:ds1:tmap:lips}) showed
activation in the ventrolateral motor cortex, corresponding to the
face/lips region of the homunculus. Additional activity was seen in
the supplementary motor area and in auditory cortex, which may
reflect the self-monitoring of facial movement.

\begin{figure*}[!t]
  \centering
  \includegraphics[width=\textwidth]{figures/fig15_ds1_tmap_lips_pos}
  \caption{Thresholded t-map for Lips$\,{>}\,$Baseline. The peak
    activation is in the ventrolateral primary motor cortex (lips
    region of the homunculus), with secondary clusters in the
    supplementary motor area and adjacent premotor cortex.}
  \label{fig:ds1:tmap:lips}
\end{figure*}

\textbf{All-motor omnibus contrast.}
The AllMotor$\,{>}\,$Baseline contrast (Fig.~\ref{fig:ds1:tmap:allmotor})
pools all three conditions and highlights the entire motor strip
together with the cerebellum and thalamus. This map gives a useful
overview of the motor system engaged by the task as a whole.

\begin{figure*}[!t]
  \centering
  \includegraphics[width=\textwidth]{figures/fig16_ds1_tmap_allmotor_pos}
  \caption{All-motor omnibus contrast (average of finger, foot, and
    lips regressors, each weighted $1/3$) for Dataset~1. The
    activation spans the full extent of primary motor cortex from
    foot (medial) to face (lateral) territory, plus the cerebellum
    and basal ganglia — a broad motor network consistent with
    voluntary movement across multiple body parts.}
  \label{fig:ds1:tmap:allmotor}
\end{figure*}

% -----------------------------------------------------------------------
\subsection{GLM Results --- Dataset 2 (Visual Task)}

The experimental timing for Dataset~2 is shown in
Fig.~\ref{fig:ds2:design_timing} and the resulting design matrix in
Fig.~\ref{fig:ds2:design_matrix}. The event-related structure of the
visual paradigm produces a denser, less obviously block-like
regressor pattern, which is visible in the design matrix as
overlapping columns.

\textbf{Face and object activations.}
Figure~\ref{fig:ds2:tmap:face} shows the face$\,{>}\,$Baseline
contrast, which produced a right-lateralised cluster in the fusiform
gyrus — the expected location of the fusiform face area (FFA),
first identified by Kanwisher et al. This replication of a classic
finding provides a useful sanity check on the pipeline.

\begin{figure*}[!t]
  \centering
  \includegraphics[width=\textwidth]{figures/fig19_ds2_tmap_face_pos}
  \caption{Thresholded t-map for faces$\,{>}\,$Baseline in
    Dataset~2. The dominant cluster in the right fusiform gyrus
    corresponds to the fusiform face area (FFA), a canonical
    face-selective region. Additional occipital activation reflects
    low-level visual processing shared across categories.}
  \label{fig:ds2:tmap:face}
\end{figure*}

\begin{figure*}[!t]
  \centering
  \includegraphics[width=\textwidth]{figures/fig21_ds2_tmap_object_pos}
  \caption{Thresholded t-map for one of the object categories
    (bottles/objects)$\,{>}\,$Baseline in Dataset~2. Activation is
    distributed across the lateral occipito-temporal cortex, more
    bilateral than the face map and extending into the fusiform
    body area (FBA). The pattern is consistent with previous
    reports from the same dataset~\cite{Haxby2001}.}
  \label{fig:ds2:tmap:object}
\end{figure*}

\begin{figure*}[!t]
  \centering
  \includegraphics[width=\textwidth]{figures/fig22_ds2_tmap_allvisual_pos}
  \caption{All-visual omnibus contrast for Dataset~2 (mean of all
    visual category regressors). The entire ventral visual stream
    from early occipital cortex through the fusiform and parahippocampal
    gyri is recruited, reflecting the common visual processing demands
    shared across all object categories.}
  \label{fig:ds2:tmap:allvisual}
\end{figure*}

% -----------------------------------------------------------------------
\subsection{Functional Connectivity}

\textbf{Dataset~1 --- Motor Network.}
Figure~\ref{fig:ds1:connectivity} shows the connectivity map seeded
from left M1. Regions correlated with the seed (warm colours,
$r>0.30$) include bilateral primary motor cortex, the supplementary
motor area (SMA) on the medial wall, premotor cortex, and the
thalamus. This spatial pattern closely matches the classic motor
resting-state network identified by Biswal et al.~\cite{Biswal1995}
in their original seed-based connectivity study. Regions showing
negative correlation (cool colours) include the default mode network
— an expected anti-correlation, since motor execution typically
suppresses default-mode activity.

\begin{figure}[!htbp]
  \centering
  \includegraphics[width=\columnwidth]{figures/fig21_ds1_connectivity_M1_hand}
  \caption{Seed-based functional connectivity map for Dataset~1.
    Seed: left primary motor cortex at MNI $[-38, -26, 60]$\,mm
    (6\,mm radius sphere). Warm voxels ($r>0.30$) form the
    sensorimotor network (bilateral M1, SMA, premotor). Cool voxels
    show anticorrelated default-mode regions.}
  \label{fig:ds1:connectivity}
\end{figure}

\textbf{Dataset~2 --- Visual Network.}
The V1 seed connectivity map (Fig.~\ref{fig:ds2:connectivity}) shows
strong positive correlations throughout the occipital lobe, extending
into the ventral and dorsal visual streams. The fusiform gyrus and
lingual gyrus show particularly high correlation, reflecting shared
responses to visual stimulation across the run. Negative
correlations are again found in medial prefrontal and posterior
cingulate regions — the default mode network — consistent with
suppression during visual processing.

\begin{figure}[!htbp]
  \centering
  \includegraphics[width=\columnwidth]{figures/fig22_ds2_connectivity_V1_calcarine}
  \caption{Seed-based connectivity map for Dataset~2. Seed: primary
    visual cortex at MNI $[0, -88, 2]$\,mm (6\,mm radius). Warm
    voxels span the occipital lobe, fusiform gyrus, and dorsal
    parietal cortex — the visual network. Cool voxels reflect
    suppression of the default mode.}
  \label{fig:ds2:connectivity}
\end{figure}

\textbf{Cross-dataset connectivity.}
Figure~\ref{fig:comparison:connectivity} puts both connectivity maps
side by side at the same axial slice. The spatial specificity is
clear: the motor seed produces connectivity in central and frontal
regions, while the visual seed is tightly constrained to occipital
and temporal areas. This double dissociation confirms that the
connectivity analysis is reflecting genuine functional network
organisation rather than a generic haemodynamic or noise effect.

\begin{figure*}[!t]
  \centering
  \includegraphics[width=\textwidth]{figures/fig23_ds1_vs_ds2_connectivity_comparison}
  \caption{Side-by-side connectivity maps: Dataset~1 motor network
    (top) and Dataset~2 visual network (bottom), at a shared axial
    slice. The spatial separation of warm-coloured clusters between
    the two maps confirms system-level functional specificity. Colour
    scale: warm = positive correlation ($r>0.30$), cool = negative.}
  \label{fig:comparison:connectivity}
\end{figure*}

% -----------------------------------------------------------------------
\subsection{Cross-Dataset Activation Comparison}

Figure~\ref{fig:comparison:activation} shows t-maps from the first
contrast of each dataset across four axial slices. The spatial
separation is immediately obvious: Dataset~1 activations sit near the
top of the brain (motor strip), while Dataset~2 activations are
concentrated in the posterior and inferior portions (occipital and
temporal cortex). This double dissociation, obtained with identical
preprocessing and GLM settings on both datasets, gives us confidence
that the pipeline produces anatomically meaningful results rather than
spurious, diffuse activations.

\begin{figure*}[!t]
  \centering
  \includegraphics[width=\textwidth]{figures/fig26_ds1_vs_ds2_activation_comparison}
  \caption{Cross-dataset activation comparison. Top row: Dataset~1
    (motor task), first contrast (finger or foot tapping). Bottom
    row: Dataset~2 (visual task), first contrast (face or object
    category). Four axial slices are shown for each. The spatial
    separation of activations — motor strip vs.\ occipito-temporal
    cortex — validates the functional specificity of the GLM results.}
  \label{fig:comparison:activation}
\end{figure*}


% ==========================================================================
\section{Discussion}
\label{sec:discussion}

The results across both datasets were broadly consistent with
expectations from the neuroimaging literature, which suggests the
pipeline is working correctly.

\textbf{Motion quality.}
Mean framewise displacement was below 0.2\,mm in both scans, putting
them in the ``low motion'' category under the criteria of Power
et~al.~\cite{Power2012}. This matters because even small amounts of
systematic motion — particularly correlated with the task — can
produce spurious activations that mimic genuine BOLD responses. The
clean motion profiles here make it unlikely that the activation maps
are motion artefacts.

\textbf{Motor cortex somatotopy.}
The three finger/foot/lips contrasts in Dataset~1 produced activation
peaks that followed the expected somatotopic organisation of primary
motor cortex from medial (foot) to lateral (face/lips)
territory~\cite{Lotze2000}. The omnibus AllMotor map covered the
entire motor strip plus the cerebellum and thalamus — the core nodes
of the voluntary motor system. These results validate both the
preprocessing (correct spatial normalisation) and the GLM specification
(correct event timing).

\textbf{Visual category selectivity.}
The face-selective fusiform cluster in Dataset~2 (Fig.~\ref{fig:ds2:tmap:face})
matches the reported location of the fusiform face area in the
original Haxby et~al.\ paper~\cite{Haxby2001}. Replicating a
published finding using an independent implementation of the same
pipeline is a strong form of validation. The other category contrasts
produced distributed activation in the lateral occipito-temporal
cortex, consistent with object-selective processing in the ventral
stream.

\textbf{Functional connectivity.}
Both connectivity maps showed spatially coherent networks. The motor
seed produced bilateral M1 and SMA connectivity (the sensorimotor
network), and the visual seed produced occipito-temporal connectivity
(the visual network). The cross-dataset comparison in
Fig.~\ref{fig:comparison:connectivity} shows these two networks are
spatially non-overlapping, which is what we would expect from
independent functional systems. The negative correlations with
default-mode regions in both maps are also expected: task-positive
networks tend to be anti-correlated with the default mode
network~\cite{Biswal1995}.

\textbf{Limitations.}
This analysis covers a single subject from each dataset, so the
results cannot be generalised beyond these two individuals. The
connectivity analysis was performed on the full task run (not a
separate resting-state scan), so the correlation structure partially
reflects task-related coactivation rather than purely intrinsic
connectivity. Future work should include group-level analyses and
separate resting-state acquisitions.

% ==========================================================================
\section{Conclusion}
\label{sec:conclusion}
% conclusion.tex
This project set out to build a complete, reproducible fMRI pipeline
in MATLAB and SPM25, and to test it on two structurally different
task paradigms. Looking at the results as a whole, the pipeline
delivered what it was supposed to.

Preprocessing quality was good across both datasets. Motion was low,
coregistration overlays showed accurate alignment, and tSNR values
were in the normal range for 3\,T acquisitions. The GLM analyses
produced spatially sensible activation maps: the three-condition
somatotopic motor map for Dataset~1 matched the known homuncular
organisation of primary motor cortex almost exactly, and the
category-selective visual responses in Dataset~2 — particularly the
face-selective fusiform cluster — replicated a well-known result from
the original Haxby et~al.\ study~\cite{Haxby2001}. Functional
connectivity analyses revealed coherent system-level networks anchored
to the seed regions, and the cross-dataset comparison confirmed that
the two datasets were comparable in data quality while differing
appropriately in their functional topography.

From an engineering perspective, the key design decision was to keep
all processing scripted and sequential (seven numbered MATLAB files),
with skip-if-exists logic so that individual steps can be re-run or
debugged without reprocessing the entire pipeline. Every intermediate
output is saved with a consistent naming scheme, making it easy to
identify which file came from which processing stage.

There is clearly room to extend this work. The analysis here covers
a single subject per dataset; a proper group study would require
second-level random-effects models and family-wise or FDR correction
for multiple comparisons. The connectivity analysis could be made
more rigorous by regressing out global signal and white matter / CSF
noise before computing correlations. Dynamic connectivity approaches
could reveal how network structure evolves across the scan. But for
the scope of this course project, the pipeline achieves its goals:
it works, it is reproducible, and the results make neurobiological
sense.


% ==========================================================================
\section*{Acknowledgement}
The author thanks Dr.\ Adnan Shah (FCSE, GIKI) for supervision and
guidance throughout this project. The OpenNeuro platform and the
contributors of datasets ds000114 and ds000105 are gratefully
acknowledged.

% ==========================================================================
\bibliographystyle{IEEEtran}
\bibliography{bibliography}

\end{document}
